\section{Sources and Approach}

The evidence base for this paper is deliberately compact and transparent. The curated corpus now contains twenty-four records:

\begin{itemize}[leftmargin=1.4em]
  \item twenty-one primary or near-primary documents from \emph{Founders Online}, the \emph{Papers of the War Department}, and Jefferson's 1808 presidential message on the Wilkinson inquiry;
  \item three local or public-history sources tied specifically to Frederick and Wilkinson's Maryland geography.
\end{itemize}

The package now releases this corpus in a more explicit audit form as well: a compact inventory with stable URLs and repository identifiers, plus a traceability file preserving the exact excerpts and summaries used for coding.

The primary records focus on three clusters. The first cluster covers early access and usefulness: the 1800 Frederick Town receipt for Wilkinson's suite, Jefferson's 1801 effort to reach both Wilkinson and Meriwether Lewis through the same forwarding network, Wilkinson's 1802 gift of river-water samples and fossils, the 1804 dinner invitation, and the 1805 buffalo-hide map letter \citep{wardepartment1800receipt,jefferson1801,wilkinson1802,wilkinson1804jefferson,wilkinson1805}. The second cluster centers on the Burr crisis, where Jefferson's support and Wilkinson's anxious replies reveal a relationship of mutual need rather than sentiment. It now also includes Jefferson's January 1808 special message, which shows him moving the Wilkinson matter into a court of inquiry while still managing it as an executive problem \citep{jefferson1807,wilkinson1807a,wilkinson1807b,jefferson1808message}. The third cluster covers rehabilitation and afterlife. It includes Wilkinson's 1811 memoir campaign and Jefferson's reassuring response, but it now also includes a fuller Frederick procedural chain: venue planning inside the War Department, Madison's refusal to preempt the Frederick tribunal, witness orders, interrogatories sent from Frederick Town, Wilkinson's own complaint to Madison, Adams's testimonial reply, Madison's approval of the acquittal, and the 1818 exchange of retrospective esteem \citep{wilkinson1811,jefferson1811,eustis1811a,madison1811eustis,eustis1811b,kinson1811,wilkinson1811madison,adams1811wilkinson,madison1812,jefferson1818,wilkinson1818}.

The local sources were chosen for function, not decoration. The Frederick Barracks marker preserves the claim that the barracks served both Lewis and Clark and later War of 1812 military traffic \citep{frederickbarracks}. The Clements Library finding aid preserves the biographical note that Wilkinson practiced medicine in Monocacy, Maryland \citep{clementsfindingaid}. The Maryland Historical Magazine article on Roger Brooke Taney preserves the later tradition that Frederick lawyers defended Wilkinson without fee during the court-martial \citep{delaplaine1918}. Read against the older Taney biographies, the Thomas Johnson and John Hanson Thomas biographical materials at the Maryland State Archives, and the \textit{Frederick-Town Herald} directory entry at the Library of Congress, that tradition belongs to a real Frederick legal and print world even if it does not settle every factual detail on its own \citep{swisher1935,smith1936,huebner2010,thomasbio,thomasjohnsonbio,fredericktownheraldloc}. None of these sources is sufficient alone; together they thicken the local setting around a primary Frederick dossier that is now much stronger than it was at the start of this project.

The claim structure is therefore tiered. The paper's strongest argument --- that Jefferson and Wilkinson were linked more by utility, protection, and political management than by simple intimacy --- rests mainly on the eleven direct Jefferson--Wilkinson records and the surrounding official correspondence. The Frederick argument is now strongest where Frederick appears in primary or near-primary documents as a venue of reputational defense, court-martial procedure, evidence gathering, and presidentially acknowledged acquittal. It is more interpretive where the local role depends on later public-history or biographical framing. Making that distinction explicit matters for a local-history audience: Frederick is important here, but not every part of its importance stands on the same evidentiary footing.

The interpretive frame, however, is broader than the corpus. This reading of Wilkinson is informed by older and newer biography, by scholarship on the Burr crisis, by studies of honor, rumor, and political reputation in the early republic, and by work on military justice and administration \citep{jacobs1938,linklater2009,cox2023,lewis2017,freeman2002,kierner2024,lurie2001,kohn1975,stagg1983,watson2012}. Those works do not supply the paper's local evidence. They do help explain why the language of utility, defense, esteem, venue, and procedural settlement belongs to the period rather than to retrospective editorial invention.

\begin{table}[t]
\centering
\small
\begin{tabular}{p{0.58\linewidth}r}
\toprule
Evidence tier & Records \\
\midrule
Direct Jefferson--Wilkinson primary correspondence & 11 \\
Primary Frederick venue and outcome records & 3 \\
Contextual Frederick local sources & 3 \\
Primary share of Frederick-linked corpus & 50\% \\
\bottomrule
\end{tabular}
\caption{Evidence ladder for the paper's core claims. The Jefferson--Wilkinson thesis rests mostly on direct correspondence, while the Frederick synthesis combines a three-document primary spine with three contextual local sources.}
\label{tab:evidence}
\end{table}


Table~\ref{tab:evidence} makes that hierarchy visible. The Frederick claim is not built from one sort of record alone. It now stands on a nine-document primary chain --- from the 1800 Frederick Town receipt and the 1804 Frederick Town letter to Adams through the 1811 War Department and interrogatory record to Madison's 1812 approval of the Frederick-town acquittal --- then thickens that chain with three local contextual sources. That layered structure is exactly why the paper can make a confident local argument without pretending every sentence has the same documentary force.

\begin{table}[t]
\centering
\small
\begin{tabular}{p{0.08\linewidth}>{\raggedright\arraybackslash}p{0.28\linewidth}>{\raggedright\arraybackslash}p{0.51\linewidth}}
\toprule
Year & Document & Function in argument \\
\midrule
1804 & Wilkinson to Adams on fame & Shows Frederick as an early site of reputational self-defense \\
1811 & Eustis on venue & Shows Frederick chosen in advance as the intended court-martial venue \\
1811 & Madison on jurisdiction & Shows the president leaving the Frederick tribunal in possession of the case \\
1811 & Eustis on witness orders & Shows the War Department mobilizing officers for the Frederick proceedings \\
1811 & Wilkinson on fairness & Shows Frederick as the place where Wilkinson sought procedural fairness \\
1811 & Kinson with interrogatories & Shows Frederick reaching outward to collect testimony by interrogatory \\
1812 & Madison on acquittal & Puts Frederick into the federal record through approved acquittal \\
1811 & Adams in reply & Shows national testimony flowing back into the Frederick defense record \\
\bottomrule
\end{tabular}
\caption{The Frederick primary documentary chain used in the paper's core local argument. These records show venue choice, jurisdiction, evidence gathering, plea, and settlement before later contextual sources are added.}
\label{tab:frederickspine}
\end{table}


Table~\ref{tab:frederickspine} makes the Frederick dossier inspectable in a second way: not just by count, but by function. The primary records do different argumentative jobs. Taken together, they show venue choice, executive restraint, witness management, outward evidence collection, plea, testimonial return, and formal federal settlement.

The coded layer is modest by design. It assigns each record a phase, relationship mode, and Frederick role; generates an event timeline; counts a small lexicon of warmth, utility, and defense terms in the excerpted correspondence; and builds a weighted network from the named relationships in the corpus. These outputs are heuristic. They do not decide the argument on their own, and network position should not be confused with causal proof. Their value is disciplinary rather than magical. If the prose claims that the relationship was mostly instrumental, or that Frederick was more than incidental, the documentary pattern has to carry that weight. The released inventory, traceability, sensitivity, and environment files are intended to make that weight easier to inspect.

The method also makes the paper easier to challenge in concrete ways. A corpus dominated by private-social records would weaken the claim that usefulness mattered more than friendship. A Frederick set made up only of late local memory would weaken the argument that Frederick was procedurally important rather than ceremonially convenient. The evidence is laid out this way so that readers can test those possibilities directly instead of taking the interpretation on trust.

\begin{table}[t]
\centering
\small
\begin{tabular}{p{0.52\linewidth}r}
\toprule
Metric & Value \\
\midrule
Coded records in corpus & 24 \\
Frederick-linked records & 12 (50\%) \\
Frederick-linked records in rehabilitation phase & 7 \\
Non-social survival modes & 18 \\
Pure social-access mode & 1 \\
Retrospective-esteem mode & 2 \\
\bottomrule
\end{tabular}
\caption{Core metrics from the coded corpus. The counts support a relationship defined much more by use and protection than by simple friendship.}
\label{tab:keyfindings}
\end{table}


The first descriptive result is already revealing. Of the twenty-four coded records, twelve are Frederick-linked. Only one record is primarily a social-access document, the 1804 dinner invitation note. By contrast, eighteen turn on use, defense, procedure, testimony, or settlement. The Frederick-linked subset is also not a one-note cluster: it spans logistical presence, reputational repair, venue choice, jurisdiction, witness coordination, interrogatory collection, trial venue, trial outcome, external testimony, infrastructure, local origin, and legal alliance. The numbers do not replace reading, but they do caution against an overly sentimental portrait.
